\section{How the Two Models Connect}
\label{sec:connect}

CLEWS and OG-Core are each detailed where the other is not. CLEWS chooses the
technologies, infrastructure and resource use that meet a given demand at the lowest cost,
and reports the resulting costs, investment and emissions. The demand it must meet is set
in advance. OG-Core determines prices, wages, interest rates, saving, investment, taxes and
public debt as households, firms and government respond to one another over time, but it
has no representation of energy technologies. What passes between the two models therefore
differs by direction. From the resource plan come costs, investment requirements and
emissions. From the economy come demand and the return on capital
(Figure~\ref{fig:connect}).

\begin{figure}[H]
\centering
\begin{tikzpicture}[
  model/.style={draw, rounded corners=2pt, thick, minimum width=4.3cm, minimum height=2.7cm,
                align=center, font=\small},
  flow/.style={-{Stealth[length=2.4mm]}, thick},
  lab/.style={font=\footnotesize, align=center}]
  \node[model, draw=clews, fill=clews!6] (C) at (0,0)
    {\textbf{CLEWS}\\[3pt] what to build and\\ how to operate it\\[3pt]
     {\footnotesize technologies, infrastructure,}\\ {\footnotesize costs, emissions}};
  \node[model, draw=ogc, fill=ogc!6] (O) at (10.4,0)
    {\textbf{OG-Core}\\[3pt] how the economy\\ responds\\[3pt]
     {\footnotesize households by age and income,}\\ {\footnotesize industries, government}};
  \draw[flow, clews] ([yshift=6mm]C.east) --
    node[lab, above] {electricity costs $\cdot$ public infrastructure\\
                      generation technology $\cdot$ air pollution} ([yshift=6mm]O.west);
  \draw[flow, ogc] ([yshift=-6mm]O.west) --
    node[lab, below] {demand for electricity\\ cost of capital} ([yshift=-6mm]C.east);
  \node[draw=policy, rounded corners=2pt, fill=policy!5, font=\footnotesize, align=center]
    (P) at (5.2,2.7) {\textbf{Policy}\quad carbon price $\cdot$ investment incentives};
  \draw[flow, policy] (P.west) -| (C.north);
  \draw[flow, policy] (P.east) -| (O.north);
\end{tikzpicture}
\caption{How OG--CLEWS connects the two models. The resource plan passes costs, investment
and pollution to the economy; the economy returns its demand and cost of capital; policy
instruments can act on both.}
\label{fig:connect}
\end{figure}

Eight channels carry these exchanges (Table~\ref{tab:channels}). Four carry outcomes of
the resource plan into the economy, two introduce policy instruments, and two return
information from the economy to the resource plan. Each channel answers a distinct
question. A scenario uses the channels relevant to the policy being assessed, so a channel
can be part of the framework without being active in every scenario.

\begin{table}[H]
\centering\small
\caption{The eight channels of OG--CLEWS and the question each answers.}
\label{tab:channels}
\begin{tabularx}{\linewidth}{@{}>{\bfseries\raggedright\arraybackslash}p{0.27\linewidth} L@{}}
\toprule
Channel & Question \\
\midrule
\multicolumn{2}{@{}l}{\emph{Resource outcomes entering the economy}} \\
Electricity cost & How do changes in the cost of electricity affect households and industries? \\
Public infrastructure & How does the infrastructure the plan requires affect production and public finances? \\
Generation capital intensity & How does a different technology mix change the way electricity is produced? \\
Air pollution and health & How do changes in air pollution affect people and the economy? \\
\addlinespace
\multicolumn{2}{@{}l}{\emph{Policy instruments}} \\
Carbon price & How does a carbon price change technology choices, payments and public revenue? \\
Investment incentives & How does policy change the incentive to invest in generation? \\
\addlinespace
\multicolumn{2}{@{}l}{\emph{Economic feedback to the resource plan}} \\
Planning discount rate & How should the economy's cost of capital inform the plan's investment choices? \\
Demand & How does the economy's response change what the resource system must supply? \\
\bottomrule
\end{tabularx}
\end{table}
